\documentclass[pdflatex,sn-nature]{sn-jnl}

\usepackage{graphicx}
\usepackage{amsmath,amssymb}
\usepackage{booktabs}
\usepackage[T1]{fontenc}
\usepackage{lmodern}
\usepackage{textcomp}
\usepackage{pdflscape}

\raggedbottom
\makeatletter
\newcommand\tabinput[1]{\@@input #1 }
\makeatother
\renewcommand{\tablename}{Supplementary Table}
\renewcommand{\thetable}{\arabic{table}}

\begin{document}

\title[Supplementary Information]{Supplementary Information for: Decoder decay and its correction over years in
intracortical brain--computer interfaces}

\author*[1]{\fnm{Md. Imtiaj Alam} \sur{Sajin}}\email{imtiajsajin@gmail.com}
\author[1]{\fnm{Esm E Moula} \sur{Chowdhury Abha}}\email{esmechowdhuryabha@gmail.com}
\author[1]{\fnm{Md Wahiduzzaman} \sur{Suva}}\email{wshuvo360@gmail.com}
\affil*[1]{\orgdiv{Department of Computer Science}, \orgname{American International University-Bangladesh},
\orgaddress{\city{Dhaka}, \postcode{1229}, \country{Bangladesh}}}

\abstract{This document contains Supplementary Notes 1--3 and Supplementary Tables 1--13. Every table is generated
from the released result files by \texttt{scripts/make\_supp\_tables.py}, so its values match the code repository.}

\maketitle

\section*{Supplementary Note 1: Label-free monitoring statistics}

Before the main analyses, we asked whether label-free statistics computed on a later day can predict how well a
renormalized fixed decoder will perform on that day. The pilot used the first 46--61 sessions of monkey N
(January 2020 to about March 2021) and 283 session pairs at most 60 days apart. Decoders were ridge regressions on 8
lags of spike-band power ($\alpha = 0.1$), trained on the first 300 trials and tested on the rest.

We computed 17 statistics in nine groups that need no labels on the later day:
\begin{itemize}
\item the shift in channel means and in channel standard deviations;
\item the principal angles between the dominant subspaces of the two days;
\item disagreement within a bootstrap ensemble of decoders;
\item disagreement between decoders built on different views of the data (spike-band power versus threshold
crossings, and 8 lags versus 1 lag);
\item statistics of the predicted output (velocity bias, variance ratio, Wasserstein distance to the training
kinematics and smoothness);
\item self-consistency between predicted position and predicted velocity;
\item Procrustes alignment disagreement, the Procrustes residual and the share of decoder readout outside the
dominant subspace;
\item neural modulation depth;
\item the log ratio of impedances and the number of silent channels.
\end{itemize}

Elapsed time predicted decoder accuracy well (Spearman $\rho = -0.75$ across all pairs and $-0.56$ within 60 days).
After controlling for log elapsed days, the partial Spearman correlation of every statistic with $R^2$ was at most
0.27 (output variance) and mostly below 0.1. In time-blocked cross-validated prediction of $R^2$, elapsed days alone
gave a mean absolute error of 0.064. The label-free statistics alone gave 0.067--0.071, and elapsed days together with
the statistics gave 0.067--0.070. When the target was the gain from recalibration rather than $R^2$, elapsed days gave
0.045 and the statistics gave 0.054. For detecting unexpectedly poor days (residual more than one standard deviation
below the time trend), every statistic had an area under the receiver operating characteristic curve of 0.50--0.58.

The residual variation was not noise. Split-half reliability of the residuals, computed on alternating pairs of
trials, was $r = 0.72$, or 0.84 after Spearman--Brown correction. Alternating pairs were needed because center-out
trials alternate between outward and return targets, which made odd/even single-trial splits negatively correlated.
Day-to-day performance beyond the time trend is therefore real and reliable, but none of the label-free statistics we
tried captured it. Abrupt faults on single channels remain detectable by construction, and such faults may still
justify monitoring.

\section*{Supplementary Note 2: Simulator}

The simulator (\texttt{ibci.sim}) is a feature-level generative model of chronic change. It transforms a real base
session into a simulated session $\Delta t$ days later, so that behavior and latent dynamics stay realistic. Its
components are defined in the Methods and its parameters are listed in Supplementary Table~12. Two further components,
heavy-tailed gain jumps and random walks of raw offsets and scales, are included for decoders that are not
renormalized. They do not affect renormalized decoders and were not calibrated.

\textit{Calibration.} The simulated ladder must be scored with exactly the same correction code and regularization as
the real ladder. An earlier calibration used regularization 100--1,000 times stronger than the values chosen on real
data. That handicapped the simulated corrections, and the optimizer compensated with distorted parameters (loss 0.663
versus 0.074 after the fix). We also restricted silencing to channels that were active in the base session, because
most channels are already silent on any given day.

\textit{Validation.} The simulator calibrated on days before 700 was evaluated on base sessions after day 700, for
all rungs of the ladder and for 20, 50 and 100 labeled trials. Two references were used: the simulator with all drift
switched off, and the median retention measured on real pairs before day 700 (an empirical curve). The simulator
beat the no-drift reference everywhere. It beat the empirical curve only for corrections it was not fitted to; for
fitted corrections and for different amounts of labeled data the empirical curve was as good or better (main text,
Fig.~8b). The simulator is therefore most useful for predicting how untested corrections behave.

\textit{Intended uses and limits.} The simulator is meant for stress-testing decoders and label-free methods under
controlled drift and electrode switching, and for comparing recalibration policies. It is calibrated on one monkey
with Utah arrays, spike-band power and a finger task, and it has not been calibrated to human recordings. Its mechanisms reproduce the measured statistics, which does not show that they are the true
physical mechanisms. Users can calibrate it to other datasets with \texttt{scripts/calibrate\_sim.py}.

\section*{Supplementary Note 3: Analysis choices revised during the study}

We report analysis choices that we revised after internal review, because each could mislead similar studies.

\textit{One-day pairs included two-day gaps.} Session pairs were matched to a target gap within 25\% or one day,
whichever was larger, so the one-day set contains two-day gaps. In participant T6 most ``one-day'' pairs were two-day
pairs, which changed the apparent overnight retention. We now report one- and two-day pairs separately.

\textit{Decoder regularization confounds decay.} With the small ridge penalties used in the original dataset papers,
decoders decayed faster than with a tuned penalty, especially in monkey N. Decay is therefore reported for decoders whose penalty
was tuned.

\textit{Label-free alignment was compared with renormalization.} In our ladder, alignment was first applied to
renormalized features, which cannot separate the two. We therefore also aligned mean-updated features.

\textit{Electrode revival needs a null.} Counting electrodes that fall silent and later recover suggested a striking
recovery rate. The same statistic computed on shuffled session orders was almost as high, so only long silences carry
evidence of a temporal process. Electrode-level tests also treat the array, not the electrode, as the unit, because
neighboring electrodes are correlated.

\textit{Flexible corrections need scrambled-channel controls.} Re-weighting one weight per channel recovered most
of the loss in the human participants, which at first suggested that drift there is mostly per-channel. Decoders
with scrambled channels could also be re-weighted to high accuracy when the decoded variable was two-dimensional,
so high recovery alone does not show per-channel drift. We therefore report scrambled-channel and non-negative
controls (Supplementary Table~5).

\textit{A full input remap is not a structural test.} A free $C\times C$ linear map in front of a frozen decoder can
reproduce any decoder of the same form when the readout rank is below the number of channels. We used it only as a
calibration target for the simulator.

\textit{Recurrent networks memorize closed-loop sessions.} A 256-unit network trained without early stopping reached
training $R^2$ of 0.97 and test $R^2$ of $-1.1$ on a human session. All network results use a smaller, regularized
network with early stopping on held-out training trials.

\textit{Small-budget recalibration shrinks the intercept.} Refitting the intercept freely with 10--20 trials shifted
the output offset and made recalibration harmful in many pairs.

\clearpage

\begin{landscape}
\begin{table}[p]
\caption{\textbf{Datasets.} Session counts are per analysis; some human sessions contribute to more than one analysis.
SBP, spike-band power; TC, threshold crossings.}
\label{tab:data}
\scriptsize\setlength{\tabcolsep}{3.5pt}
\begin{tabular}{@{}llllll@{}}
\toprule
Individual & Source & Arrays and site & Signal & Task & Sessions (span) \\
\midrule
Monkey N & DANDI 001201 (LINK) & 2 Utah, motor cortex & SBP, TC & finger flexion & 312 on 303 days (3.4 years) \\
Monkey C & DANDI 000688 & Utah, motor/premotor & sorted spikes & reaching & 68 (3.0 years) \\
Monkey M & DANDI 000688 & Utah, motor/premotor & sorted spikes & reaching & 28 (1.5 years) \\
\tabinput{supp/tab_data_humans.tex}
14 humans & Dryad x0k6djj1h & 20 Utah arrays & TC rates, impedance & research sessions & 2,289 ($\leq$7.3 years) \\
T11, T5 & Dryad n2z34tn5s (MINDFUL) & motor cortex & SBP (T11), TC (T5) & closed-loop cursor & 15 and 6 days \\
\botrule
\end{tabular}
\end{table}
\end{landscape}

\begin{landscape}
\begin{table}[p]
\caption{\textbf{Correction ladder by individual and gap (linear decoder, default penalty).} Median decoding $R^2$ for
each rung, and median retention of the renormalized fixed decoder with 95\% cluster-bootstrap confidence intervals.
Each training session contributes at most one pair per gap. Supervised rungs used up to 300 labeled trials.
Aligned, Procrustes rotation after renormalization; stable aligned, the same estimated from stable channels.}
\label{tab:ladder}
\scriptsize\setlength{\tabcolsep}{2.8pt}
\begin{tabular}{@{}rrrrrrrrrrl@{}}
\toprule
Gap (days) & Pairs & Reference & None & Mean update & Renorm. & Aligned & Stable aligned & Re-weighted & Ridge to prior &
Retention [95\% CI] \\
\midrule
\tabinput{supp/tab_ladder.tex}
\botrule
\end{tabular}
\end{table}
\end{landscape}

\begin{landscape}
\begin{table}[p]
\caption{\textbf{Label-free corrections.} Median over training sessions of the paired difference in $R^2$ (pairs at
gaps up to 120 days for alignment, all gaps for CORAL and the stabilizer), with two-sided Wilcoxon signed-rank $P$ in
parentheses. Mean, mean updating only; rot., Procrustes rotation; stab., factor-analysis stabilizer; FA, factor
analysis.}
\label{tab:labelfree}
\scriptsize\setlength{\tabcolsep}{3.5pt}
\begin{tabular}{@{}lllllll@{}}
\toprule
Individual & Renorm. vs mean & Mean+rot. vs mean & Renorm.+rot. vs renorm. & CORAL vs renorm. & Stab. vs renorm. &
Stab. vs fixed FA \\
\midrule
\tabinput{supp/tab_labelfree.tex}
\botrule
\end{tabular}
\end{table}
\end{landscape}

\begin{table}[p]
\caption{\textbf{Monitoring.} \textbf{a}, Offline, monkey N (537 pairs): Spearman correlation with decoder $R^2$, raw and
partial given log elapsed days, and time-blocked cross-validated error of predicted $R^2$ (days and divergence uses
both divergences). \textbf{b}, Closed loop, public MINDFUL data: Pearson correlation of angular error with elapsed days
and each divergence (raw and partial given days), and leave-one-day-out mean absolute error (degrees).}
\label{tab:mindful}
\scriptsize\setlength{\tabcolsep}{3.5pt}
\textbf{a}\\[2pt]
\begin{tabular}{@{}lrr@{}}
\toprule
Predictor & Raw $\rho$ & Partial $\rho$ \\
\midrule
\tabinput{supp/tab_mindful_link.tex}
\botrule
\end{tabular}

\vspace{6pt}
\textbf{b}\\[2pt]
\begin{tabular}{@{}lrrrrrrrrrrr@{}}
\toprule
& & & & \multicolumn{2}{c}{Neural} & \multicolumn{2}{c}{Output} & \multicolumn{4}{c}{Error (\textdegree)} \\
\cmidrule{5-6}\cmidrule{7-8}\cmidrule{9-12}
& Days & Windows & $r$(days) & raw & partial & raw & partial & days & +neural & +output & output \\
\midrule
\tabinput{supp/tab_mindful.tex}
\botrule
\end{tabular}
\end{table}

\begin{landscape}
\begin{table}[p]
\caption{\textbf{Channel controls and electrode link.} For 40 training sessions per individual (tuned penalty), median
retention of the renormalized fixed decoder, after re-weighting the correctly assigned decoder, after re-weighting a
decoder with scrambled channels, and the same two with non-negative weights. Electrode link (humans): median over
training sessions of the Spearman correlation between fitted channel weights and the change in each electrode's log
threshold-crossing rate or log impedance between the two days, with the share of sessions with a positive correlation.}
\label{tab:controls}
\scriptsize\setlength{\tabcolsep}{3.5pt}
\begin{tabular}{@{}lrrrrrrll@{}}
\toprule
& & & \multicolumn{2}{c}{Signed weights} & \multicolumn{2}{c}{Non-negative weights} & \multicolumn{2}{c}{Electrode link} \\
\cmidrule{4-5}\cmidrule{6-7}\cmidrule{8-9}
Individual & Pairs & No labels & Correct & Scrambled & Correct & Scrambled & Firing rate & Impedance \\
\midrule
\tabinput{supp/tab_controls.tex}
\botrule
\end{tabular}
\end{table}
\end{landscape}

\begin{table}[p]
\caption{\textbf{Recurrent network versus linear decoder.} For 25 training sessions per individual: median reference
$R^2$, retention of the renormalized fixed decoder, and share of the loss recovered by re-weighting each channel.}
\label{tab:nn}
\scriptsize\setlength{\tabcolsep}{3.5pt}
\begin{tabular}{@{}rrrrrrrr@{}}
\toprule
& & \multicolumn{2}{c}{Reference $R^2$} & \multicolumn{2}{c}{Retention} & \multicolumn{2}{c}{Loss recovered} \\
\cmidrule{3-4}\cmidrule{5-6}\cmidrule{7-8}
Gap (days) & Pairs & Linear & Network & Linear & Network & Linear & Network \\
\midrule
\tabinput{supp/tab_nn.tex}
\botrule
\end{tabular}
\end{table}

\begin{table}[p]
\caption{\textbf{Matched-output control in the reaching monkeys} (tuned penalty). Median retention of the renormalized
fixed decoder and after re-weighting, and the share of the loss recovered by re-weighting (\%), when decoding hand
kinematics (four outputs) or hand-movement direction only (a two-dimensional unit vector, as for the human
participants).}
\label{tab:matched}
\scriptsize\setlength{\tabcolsep}{3.5pt}
\begin{tabular}{@{}llrrrr@{}}
\toprule
Individual & Target & Gap (days) & Renormalized & Re-weighted & Recovered (\%) \\
\midrule
\tabinput{supp/tab_matched.tex}
\botrule
\end{tabular}
\end{table}

\begin{table}[p]
\caption{\textbf{Angles between decoders fitted on different days} (degrees, medians over pairs). Global, angle that
allows one global scale; per channel, angle after fitting one scale per channel; permuted, the per-channel angle
after randomly permuting the channels of the earlier decoder.}
\label{tab:angles}
\scriptsize\setlength{\tabcolsep}{3.5pt}
\begin{tabular}{@{}lrrrr@{}}
\toprule
Individual & Gap (days) & Global & Per channel & Permuted \\
\midrule
\tabinput{supp/tab_angles.tex}
\botrule
\end{tabular}
\end{table}

\begin{landscape}
\begin{table}[p]
\caption{\textbf{Regularization by the 1\% rule.} The rule chose the largest ridge penalty whose median same-day $R^2$
was within 1\% of the best. Columns give median $R^2$ with $\alpha = 0.1$ and with the chosen $\alpha$, the median
paired gain with 95\% cluster-bootstrap confidence interval, and the share of training sessions whose mean gain was
positive.}
\label{tab:reg}
\scriptsize\setlength{\tabcolsep}{3.5pt}
\begin{tabular}{@{}llrrrrlr@{}}
\toprule
Individual & Chosen $\alpha$ & Gap (days) & Pairs & $R^2$, $\alpha = 0.1$ & $R^2$, rule & Gain [95\% CI] &
Share positive \\
\midrule
\tabinput{supp/tab_reg.tex}
\botrule
\end{tabular}
\end{table}
\end{landscape}

\begin{landscape}
\begin{table}[p]
\caption{\textbf{Recalibration with few labeled trials.} \textbf{a}, Monkey N (86 pairs): median retention without
labels and, for each number of labeled trials, ridge to prior / fresh decoder / re-weighting (regularization by
cross-validation). \textbf{b}, Human participants (30 training sessions each, all gaps): re-weighting / ridge to prior /
fresh decoder. \textbf{c}, Harmful recalibrations in monkey N (more than 0.01 $R^2$ below the renormalized fixed
decoder) with the shrinkage chosen by cross-validation or fixed at the value best on other training sessions, with
exact McNemar $P$, and median retention for both and for the best value in hindsight.}
\label{tab:recal}
\scriptsize\setlength{\tabcolsep}{3.5pt}
\textbf{a}\\[2pt]
\begin{tabular}{@{}rrrlllll@{}}
\toprule
Gap (days) & Pairs & No labels & $n = 10$ & $n = 20$ & $n = 50$ & $n = 100$ & $n = 300$ \\
\midrule
\tabinput{supp/tab_efficiency.tex}
\botrule
\end{tabular}

\vspace{6pt}
\textbf{b}\\[2pt]
\begin{tabular}{@{}lrrlllll@{}}
\toprule
Participant & Pairs & No labels & $n = 10$ & $n = 20$ & $n = 50$ & $n = 100$ & $n = 300$ \\
\midrule
\tabinput{supp/tab_eff_human.tex}
\botrule
\end{tabular}

\vspace{6pt}
\textbf{c}\\[2pt]
\begin{tabular}{@{}rrrrrrrr@{}}
\toprule
& & \multicolumn{3}{c}{Harmful (\%)} & \multicolumn{3}{c}{Median retention} \\
\cmidrule{3-5}\cmidrule{6-8}
Trials & Pairs & Cross-validation & Fixed & $P$ (Fisher) & Cross-validation & Fixed & Hindsight \\
\midrule
\tabinput{supp/tab_policy.tex}
\botrule
\end{tabular}
\end{table}
\end{landscape}

\begin{landscape}
\begin{table}[p]
\caption{\textbf{Electrode statistics.} \textbf{a}, Silence and recovery in initially active human electrodes under
different definitions, observed and after shuffling each electrode's session order (shuffled silent episodes in
parentheses). \textbf{b}, Per array: yield is the percentage of electrodes with a threshold-crossing rate of at least
2~Hz (median of the first / last three sessions); $h_{\mathrm{off}}$ and $h_{\mathrm{on}}$ are switching rates per
1,000 days; Rev./sil., recovered and silenced initially active electrodes; Kurt., excess kurtosis of
session-to-session changes in log rate after removing session-wide changes; Imp.\ $\rho$, Spearman correlation of
the array median impedance with day; Edge $-$ int., median edge minus median interior slope of log rate per year, with
the one-sided Mann--Whitney $P$ (electrodes treated as independent; see main text for array-level tests).}
\label{tab:electrodes}
\scriptsize\setlength{\tabcolsep}{3.5pt}
\textbf{a}\\[2pt]
\begin{tabular}{@{}lrrrr@{}}
\toprule
Definition & Silent episodes & Recovered & Recovered (\%) & Shuffled (\%) \\
\midrule
\tabinput{supp/tab_revival.tex}
\botrule
\end{tabular}

\vspace{6pt}
\textbf{b}\\[2pt]
\begin{tabular}{@{}llrrrrrrrrrr@{}}
\toprule
Part. & Array & Sess. & Span (d) & Yield (\%) & $h_{\mathrm{off}}$ & $h_{\mathrm{on}}$ & Rev./sil. & Kurt. &
Imp. $\rho$ & Edge $-$ int. & $P$ \\
\midrule
\tabinput{supp/tab_failure.tex}
\botrule
\end{tabular}
\end{table}
\end{landscape}

\begin{table}[p]
\caption{\textbf{Simulator parameters.} Values calibrated on monkey N sessions before day 700 (early) and from day 700
onwards (late). Switching rates were measured from electrode activity in each period. Fixed parameters: local mixing
share $\beta = 0.5$, turnover concentration 5, 20 latent dimensions and silent-signal fraction $\kappa = 0.37$.}
\label{tab:sim}
\scriptsize\setlength{\tabcolsep}{3.5pt}
\begin{tabular}{@{}llrr@{}}
\toprule
Parameter & Unit & Early & Late \\
\midrule
\tabinput{supp/tab_sim.tex}
\botrule
\end{tabular}
\end{table}

\begin{table}[p]
\caption{\textbf{Training on simulated futures.} Median $R^2$ on real later sessions of monkey N for ridge decoders
trained on sessions after day 700 (simulator calibrated before day 700). Base, $\alpha = 0.1$; stronger ridge,
$\alpha = 10^4$; perturbation, random channel dropout and gain noise; simulator, eight simulated futures; simulator and
ridge, both; several days, data from the three previous sessions; reference, decoder trained on the test day.}
\label{tab:augment}
\scriptsize\setlength{\tabcolsep}{3.5pt}
\begin{tabular}{@{}rrrrrrrrr@{}}
\toprule
Gap (days) & Pairs & Base & Stronger ridge & Perturbation & Simulator & Simulator and ridge & Several days &
Reference \\
\midrule
\tabinput{supp/tab_augment.tex}
\botrule
\end{tabular}
\end{table}

\begin{table}[p]
\caption{\textbf{BrainGate decoding participants and the inclusion rule.} Arrays (channels), sessions and span of the
closed-loop cursor sessions; median trials (movement periods) per session and median movement duration; median $R^2$ of
the reference decoder across session pairs (linear decoder, default penalty). Participants entered the decay and
correction analyses when this median exceeded 0.1 (Methods). For T2, neither block-wise normalization nor threshold
crossings raised same-day $R^2$ above about 0.1 except on its earliest session.}
\label{tab:decoding}
\small
\begin{tabular}{@{}lccccccc@{}}
\toprule
Participant & Arrays (channels) & Sessions & Span (years) & Trials per session & Movement (s) & Reference $R^2$ & Included \\
\midrule
\tabinput{supp/tab_decoding.tex}
\botrule
\end{tabular}
\end{table}

\end{document}
